@ID: ON17D1P1U
@BIDANG: Kombinatorika

Untuk bilangan bulat $k$ dan $n$ dengan $0\le k\le n$, buktikan bahwa

$$
\sum_{m=k}^{n}\binom{m}{k}
=
\binom{n+1}{k+1}.
$$

@ID: ON17D1P2U
@BIDANG: Aljabar Linear

Misalkan $A$ dan $B$ adalah matriks dalam $\mathbb R^{n\times n}$ yang memenuhi

$$
AB^2-2BAB+B^2A=0.
$$

Tentukan nilai eigen terbesar dari matriks $AB-BA$.

@ID: ON17D1P3U
@BIDANG: Analisis Kompleks

Misalkan $a$, $b$, dan $c$ adalah bilangan-bilangan kompleks dengan sifat $abc=1$ dan

$$
a^{20}+b^{20}+c^{20}
=
\frac1{a^{20}}+\frac1{b^{20}}+\frac1{c^{20}},
$$

$$
a^{17}+b^{17}+c^{17}
=
\frac1{a^{17}}+\frac1{b^{17}}+\frac1{c^{17}},
$$

$$
a^{2017}+b^{2017}+c^{2017}
=
\frac1{a^{2017}}+\frac1{b^{2017}}+\frac1{c^{2017}}.
$$

Buktikan bahwa

$$
1\in\{a,b,c\}.
$$

@ID: ON17D1P4U
@BIDANG: Analisis Real

Jika fungsi $f:\mathbb R\to\mathbb R$ mempunyai sifat bahwa untuk setiap deret $\sum_{n=1}^{\infty}a_n$ yang konvergen, deret $\sum_{n=1}^{\infty}f(a_n)$ juga konvergen, buktikan bahwa terdapat $M>0$ dan $\varepsilon>0$ sehingga untuk setiap $x\in\mathbb R$ dengan $0<x<\varepsilon$ berlaku

$$
|f(x)|\le Mx.
$$

@ID: ON17D1P5U
@BIDANG: Struktur Aljabar

(a) Berikan contoh suatu bilangan prima ganjil $p$ dan suatu grup $G$ sehingga $|G|=p+1$ dan $p$ membagi $|\operatorname{Aut}(G)|$.

(b) Misalkan $p$ suatu bilangan prima ganjil dan $G$ suatu grup dengan $|G|=p+1$. Buktikan bahwa jika $p$ membagi $|\operatorname{Aut}(G)|$, maka $p=4k+3$ untuk suatu bilangan bulat $k$.

@ID: ON17D2P1U
@BIDANG: Aljabar Linear

Misalkan $A$ dan $B$ adalah matriks dalam $\mathbb R^{2017\times2017}$ yang memenuhi

$$
A^{-1}=(A+B)^{-1}-B^{-1}
$$

dan

$$
\det(A^{-1})=2017.
$$

Tentukan $\det(B)$.

@ID: ON17D2P2U
@BIDANG: Struktur Aljabar

Misalkan $R$ suatu gelanggang yang memenuhi $x^2=x$ untuk setiap $x\in R$.

(a) Tunjukkan bahwa $R$ merupakan gelanggang komutatif.

(b) Buktikan bahwa untuk setiap $x_1,x_2,\ldots,x_n\in R$, ideal

$$
I=\langle x_1\rangle+\langle x_2\rangle+\cdots+\langle x_n\rangle
$$

dapat dituliskan sebagai $I=\langle x\rangle$ untuk suatu $x\in R$.

@ID: ON17D2P3U
@BIDANG: Kombinatorika

Sebuah barisan $\{Y_n\}$ didefinisikan oleh $Y_1=2$ dan

$$
Y_{n+1}=Y_n(Y_n-1)+1,\qquad n\ge1.
$$

Buktikan bahwa bila $m\ne n$, maka pembagi sekutu terbesar dari $Y_m$ dan $Y_n$ adalah $1$, dan

$$
\sum_{i=1}^{\infty}\frac1{Y_i}=1.
$$

@ID: ON17D2P4U
@BIDANG: Analisis Real

Diketahui fungsi $f:\mathbb R\to\mathbb R$ kontinu dan surjektif. Jika untuk setiap $y\in\mathbb R$, cacah anggota himpunan

$$
\{x\in\mathbb R:f(x)=y\}
$$

paling banyak dua, buktikan bahwa $f$ monoton.

@ID: ON17D2P5U
@BIDANG: Analisis Kompleks

Misalkan $C=\{z\in\mathbb C:|z|=2\}$ dan

$$
f(z)=z^4+\frac{2017}{z^4+z^3+1}.
$$

(a) Buktikan bahwa untuk setiap $R>2$ berlaku

$$
\int_C f(z)\,dz
=
\int_{|z|=R}f(z)\,dz.
$$

(b) Hitung

$$
\int_C f(z)\,dz.
$$